\section{连续缓存分配与换入换出}
\subsection{逻辑生命与物理驻留状态}
上一问把缓冲区生命期表示为逻辑区间。本问为这些对象分配真实连续地址，一个逻辑生命期因换出可以拆成若干段物理驻留。每次换出释放片内区间，后续使用前再恢复对象，恢复地址可以改变。因此求解同时维护“数据还需保留”和“数据当前占据哪一片地址”两种信息。

给定规范化顺序后，问题二还需要确定初始偏移、换出对象及其插入位置、配对换入位置和新偏移。问题允许进一步调整原始操作顺序，因此这些选择共同构成可行域。缓存容量分别为 $C_{L1}=4096$、$C_{UB}=1024$、$C_{L0A}=C_{L0B}=256$、$C_{L0C}=512$，类型之间不能借用空间。

每个缓冲区在任一序列位置处于四种状态之一：未申请、驻留、已换出、已释放。合法状态迁移为
\begin{equation}
\text{未申请}\xrightarrow{\rm ALLOC}\text{驻留}
\mathrel{\mathop{\rightleftarrows}^{\rm SPILL\_OUT}_{\rm SPILL\_IN}}
\text{已换出},\qquad
\text{驻留}\xrightarrow{\rm FREE}\text{已释放}.\label{eq:states}
\end{equation}
只有驻留状态允许操作引用；已释放后不可再次申请同一标识。换出不终结逻辑生命，仍须保留原始 FREE。若换出后已无普通操作需要该对象，也必须完成配对换入后再释放，以符合新增节点和配对规则，而不能删去换入以降低指标。

用序列区间 $[\pi(o_{bj}),\pi(z_{bj}))$ 表示第 $j$ 段物理驻留，从 ALLOC 或 SPILL\_IN 开始，到 FREE 或 SPILL\_OUT 结束。其地址为整数半开区间 $[x_{bj},x_{bj}+s_b)$，满足
\begin{equation}
x_{bj}\in\mathbb Z_{\ge0},\qquad x_{bj}+s_b\le C_{t_b}.\label{eq:address}
\end{equation}
若两个同类型驻留段在序列上重叠，其地址必须满足析取约束
\begin{equation}
x_{bj}+s_b\le x_{ak}\quad\text{或}\quad x_{ak}+s_a\le x_{bj}.\label{eq:disjoint}
\end{equation}
每个使用者必须落在相应驻留段内，段的关闭节点等待该段全部使用者，下一次换入等待对应换出。时间上的地址互斥还由下一节的释放到再占用依赖保证；仅在输出序列上画不重叠区间，并不足以证明并行执行时安全。

\subsection{搬运目标与连续空间的困难}
令 $\chi_b$ 按题面 COPY\_IN 引用关系取值。问题二的目标是
\begin{equation}
\min D=\sum_{(b,j)\in\mathcal M}(2-\chi_b)s_b,\label{eq:traffic}
\end{equation}
约束为原始节点拓扑顺序、式\eqref{eq:states}--\eqref{eq:disjoint}及换出换入配对。重复换出同一对象按不同事件重复计费。一次换入耗时 $2s_b+150$；换出在 $\chi_b=1$ 时耗时零，否则也是 $2s_b+150$。这些耗时进入问题三的执行单元工作量，零耗时换出仍是依赖节点。

容量和连续性是两类约束。以容量 10 为例，若两个当前需保留的对象位于 $[0,2)$、$[6,8)$，剩余空闲区间长 4 和 2；新对象长 5，三者总长仅为 9，却没有一个空闲段可容纳它。因而“空闲总量足够”不是分配成功的充分条件。动态存储分配文献将空闲区间组织与外部碎片作为核心问题\cite{wilson}，本文据此同时检查空闲总量与可用连续区间。

记某类型当前最大连续空闲段长度为 $L_t^{\max}$，空闲总量为 $F_t$。请求大小为 $s$ 时，固定现有偏移下的可分配条件是 $L_t^{\max}\ge s$；若 $F_t\ge s>L_t^{\max}$，阻塞来自外部碎片；若 $F_t<s$，则空闲总量已经不足。两种情形都可能通过换出释放空间，但牺牲对象的位置决定释放后能否合并成足够长的区间。后文的距离评分按对象选牺牲者，连续区间扫描负责确认每次牺牲是否已经达到分配条件。

最佳适配枚举所有长度不小于 $s$ 的空闲段，选择长度最小的一段，并从该段起始位置分配；长度相同再取较低地址。其直接效果是使本次分配在所选空闲段内留下的余量较小，较长空闲段可以暂时保留给后续较大对象。未来申请和释放可能改变区间结构，因此每次仍重新计算空闲段。由此，地址选择是基于当前状态的启发式，与完整调度顺序共同决定最终搬运量。

\subsection{分配、换出与恢复算法}
先沿原始顺序预存每个对象的未来使用位置队列，再逐节点模拟四状态迁移。遇到申请或换入，扫描对应缓存，优先选择长度最小且足够的空闲区间。若没有可用区间，排除当前操作必需的对象，从其余同类型驻留对象中选择牺牲者，插入换出并合并空闲区间，重复直至分配成功。

牺牲评分比较两种选择：最大化下一次使用距离 $\Delta_b$，或最大化
\begin{equation}
q_b^{\rm spill}=\frac{\Delta_b}{(2-\chi_b)s_b}.\label{eq:victim}
\end{equation}
无后续普通操作使用的对象，将其距离置为原始序列长度加一。后一规则倾向于换出未来较晚使用且搬运较便宜的对象，连续空闲段的组合效果则在每次换出后重新扫描确认。Bélády 对替换问题的研究利用完整未来引用序列构造离线比较基准\cite{belady}。本文借鉴未来引用信息，在对象大小不等、多对象同时使用及连续地址约束下构造上述两种牺牲评分，并以六图真实搬运量比较其表现。

普通换出可能被受保护对象的碎片阻塞，此时启用显式恢复：换出本次操作所需缓存类型中的全部驻留对象，再按对象大小递减顺序重装本次必需集合。该步骤允许改变偏移，但每次搬运均计入式\eqref{eq:traffic}，不是零成本地址搬移。当前实现以容量单元数组表示占用情况，便于检查边界；这也是容量不大时选择简单实现的原因。

恢复成功的前提为每个操作在每种缓存上的必需总量满足
\begin{equation}
\sum_{b\in B_v:\,t_b=t}s_b\le C_t,\qquad v\in V_o,\ t\in\mathcal T.\label{eq:necessary}
\end{equation}
这一前提已对六图逐操作检查。证明如下：将相关类型清空后，在同一类型中把必需对象紧邻排列，其最后一个地址不超过总长，而总长不超过容量，所以存在连续放置。未换入的其他对象可留在 DDR 等待后续使用；DDR 容量在题面中不作限制。在固定原始顺序下，每个节点均可通过有限次换出和重装推进，因而不会因外部碎片永久停滞。若式\eqref{eq:necessary}失败，单纯换出策略无法解决同时使用需求，应直接报告不可行，而不是无限循环。

初始偏移对应原始 ALLOC，每次换出另记录新的换入偏移。若有 $M$ 次换出，新增节点恰为 $2M$ 个，第 $j$ 次从零编号时使用 $N+2j,N+2j+1$，并保持原始节点各出现一次。未来使用队列的预处理为 $O(N+Q)$。计入容量扫描、牺牲对象扫描及恢复时的对象排序，分配阶段的保守时间界为 $O((N+M)(C+|B|+K\log(K+1))+Q)$，其中 $C=\max_tC_t$；显式恢复增加的工作已包含在实际 $M$ 内，不将其当作常数开销。

\subsection{一次分配的决策计算与恢复示例}
未来使用信息来自已经给定的原始节点顺序。设当前处理位置为 $i$，对象 $b$ 的未来使用队列首项为 $j_b$，则 $\Delta_b=j_b-i$。处理完当前操作后，弹出本次使用位置；以后即使对象暂时换出，其未来队列仍继续保留。这样，逻辑需求与物理状态相互独立，换入时能够找到下一使用需求，而无需重新扫描整个后缀。

两种牺牲规则的差别可以用一个构造状态说明。设同类型两个可换出对象 $a,b$ 的大小分别为 4、2，下一使用距离分别为 12、8；$a$ 为题面输入复制对象，$b$ 为中间对象。两者一次事件的搬运代价均为 4，距离评分和代价评分都优先换出 $a$。若把 $a$ 改为中间对象，其代价变为 8，代价评分变为 $12/8=1.5$，而 $b$ 的评分为 $8/4=2$，此时距离规则仍选择 $a$，代价规则选择 $b$。这一计算说明，分母改变的是在未来距离与本次搬运之间的取舍。

牺牲规则还受到保护集合与地址形状的共同限制。若当前操作同时使用 $a$，它就从普通牺牲集合中排除，即使其评分最大也不会被普通换出选中。若换出 $b$ 后仍无足够长区间，算法会继续选择其他合法牺牲者；一次申请因此可能对应多次换出。评分相同时，当前实现优先较大对象，再以较小对象编号打破平局，使相同输入生成相同结果。

在容量 10 的碎片小例中，当前必需的两个对象各长 2，分别位于 $[0,2)$、$[6,8)$，还需装入一个长 5 的对象。必需总量为 9，满足式\eqref{eq:necessary}，但原地址布局阻塞请求。若两对象均受保护且没有其他可牺牲对象，普通换出无法推进。显式恢复先将二者换出，再把三个必需对象按长度排列为 $[0,5)$、$[5,7)$、$[7,9)$，余下 $[9,10)$ 空闲。

若被换出的两个对象均为中间对象，该恢复至少增加 $2\times2+2\times2=8$ 的搬运量；每个对象的换出、换入各耗时 $2\times2+150=154$ 周期。这些事件分别进入输出和输入搬运单元，实际增加的完成时间还由依赖与重叠决定。该构造例展示了空间可恢复性及其代价，既解释为什么总量可行仍会遇到局部失败，也说明恢复过程需要显式记录搬运。

对每个普通操作，算法先检查各类型必需总量，然后按大小递减确保全部需求对象驻留。若部分对象已经装入而后续请求遭遇保护碎片，此前已发生的换出、换入仍保留，恢复从当前状态继续执行。该处理使搬运统计与实际生成的操作一致。输出 FREE 时若对象尚在 DDR，则先补配对换入，再释放其逻辑生命；因此每次换出最终都有对应恢复，所有驻留状态在遍历结束时归零。

恢复的有限性可由两阶段说明。普通换出每次减少一个当前驻留且未受保护对象，有限对象集合使该循环有限。清空相关类型后，大小递减装入的任一前缀总长不超过完整必需集合总长，故各次申请都可放在剩余连续后缀中。处理完本操作再进入下一节点，整个原始序列有限，从而输出完整有限方案。

\subsection{搬运量结果及策略取舍}
问题二在编号、三组缓存贪心及两种深度优先顺序上，各比较两种换出评分，共十二个候选。按 $D$ 最小选择，平局取较小执行时间。表\ref{tab:traffic}的编号基准采用相同生命周期收缩、最佳适配和距离评分，使比较反映顺序与牺牲规则的组合变化，而不是额外预处理差异。
\begin{table}[H]\centering\caption{问题二的额外搬运量与换出次数}\label{tab:traffic}\small
{\setlength{\tabcolsep}{12pt}\begin{tabular}{lrrrr}\toprule
计算图 & $D_{\rm id}$ & $D_2$ & 降幅/\% & 换出次数\\\midrule
卷积小图 & 87612 & 44766 & 48.90 & 56\\
卷积大图 & 76636 & 75546 & 1.42 & 2116\\
注意力小图 & 3620 & 3620 & 0.00 & 23\\
注意力大图 & 32852 & 32852 & 0.00 & 212\\
矩阵乘小图 & 34944 & 34944 & 0.00 & 273\\
矩阵乘大图 & 460800 & 460800 & 0.00 & 3600\\
\bottomrule\end{tabular}}
\end{table}

卷积小图的搬运量从 87612 降至 44766，降幅 48.90\%；顺序改为深度优先、仍用距离评分时已经降至 45336，再换用代价评分降至 44766。因此主要收益来自该图上的顺序变化，而代价评分进一步节省 570。卷积大图则选择深度优先与距离评分，搬运量为 75546；若对同一深度优先顺序改用代价评分，搬运量反而升至 81538。这一对照直接说明式\eqref{eq:victim}不是普遍优于距离评分的准则。

两个注意力图和两个矩阵乘图均保留编号基准的搬运量，其中矩阵乘大图有 3600 次换出、搬运量为 460800。换出次数与单次大小共同决定总量，因此下面分别考察事件分布和重复对象，以识别这一开销的具体组成。

\subsection{分类型搬运与重复换出}
对每个类型 $t$，将该类型所有换出事件的式\eqref{eq:traffic}贡献求和得到 $D_t$，满足 $D=\sum_tD_t$。再以对象标识合并事件，记发生过换出的不同对象数为 $B_s$，重复事件数定义为 $M-B_s$。这两个统计量分别反映空间类型和同一对象被反复移出的程度。表\ref{tab:spilltypes}给出问题二方案的实际结果，其他零级缓存类型的额外搬运量在六图中均为零。
\begin{table}[H]\centering\small
\caption{额外搬运的缓存类型与重复事件}\label{tab:spilltypes}
\begin{tabular}{lrrrrr}\toprule
计算图 & $D_{L1}$ & $D_{UB}$ & 不同对象 & 重复事件 & 最大次数\\\midrule
卷积小图 & 44766 & 0 & 28 & 28 & 5\\
卷积大图 & 75546 & 0 & 690 & 1426 & 9\\
注意力小图 & 0 & 3620 & 23 & 0 & 1\\
注意力大图 & 8448 & 24404 & 184 & 28 & 7\\
矩阵乘小图 & 34944 & 0 & 46 & 227 & 7\\
矩阵乘大图 & 460800 & 0 & 255 & 3345 & 15\\
\bottomrule\end{tabular}
\end{table}

两个卷积图和两个矩阵乘图的额外搬运均来自一级缓存。卷积小图有 28 个不同对象被换出，形成 56 次事件；大图有 690 个不同对象，形成 2116 次事件，单个对象最多换出 9 次。大图的重复事件数为 1426，表明搬运并非完全由大量互不相同的对象各换出一次构成。共享对象在不同需求之间被移出再恢复，其未来使用跨度与分支执行顺序因而值得共同考虑。

矩阵乘小图仅有 46 个不同对象被换出，却形成 273 次事件；大图相应为 255 个对象、3600 次事件，最大重复次数为 15。结合表\ref{tab:peaktypes}中一级缓存主导的逻辑峰值，可以把问题一的长时间保留需求与问题二的反复恢复联系起来：有限容量方案通过多次恢复服务同一批持续存在的逻辑对象。该统计定位了当前方案的开销分布，进一步降低搬运量需要重新安排对象的多次使用或选择其他驻留关系。

注意力小图的 23 次换出涉及 23 个不同对象，全部发生在统一缓存，总搬运量为 3620。注意力大图则有 179 次统一缓存换出及 33 次一级缓存换出，对应搬运量 24404 和 8448。这一差别与峰值组成相呼应：小图总峰值位置的一级缓存为 3072，低于 4096，而统一缓存为 1044，已超过 1024；大图的相应值为 5120 与 2084，两种类型都出现较强需求。分类型统计因而比单个总峰值更直接地说明换出发生的位置。

\subsection{真实驻留事件与连续地址复用}
图\ref{fig:occupancy}按扩展调度序列重放注意力小图的物理驻留。横轴包含新增换出换入节点，纵轴分别是一级缓存和统一缓存的占用量；虚线表示相应物理容量。每次申请或换入使占用上升，每次释放或换出使占用下降。此图用于观察逻辑序列中的空间状态，执行周期上的等待将在下一节通过流水图分析。
\begin{figure}[H]\centering
\includegraphics[width=.93\linewidth]{../../figures/npu2025-r3/attention_occupancy.pdf}
\caption{注意力小图问题二方案的分类型物理驻留}\label{fig:occupancy}\end{figure}

第一处换出发生在扩展序列第 237 个位置。此前统一缓存已占 908，空闲总量为 $1024-908=116$，下一个对象 87 需要 128，因此该请求首先遇到总量不足。算法将大小 128 的对象 78 从 $[520,648)$ 换出，占用降至 780，再在同一区间申请对象 87，占用回到 908。节点 1716 对应这次换出，节点 171 对应对象 87 的原始申请。具体位置与地址见表\ref{tab:events}。
\begin{table}[H]\centering\small
\caption{注意力小图首个换出附近的统一缓存事件}\label{tab:events}
\begin{tabular}{rrlrrrr}\toprule
位置&节点&操作&对象&大小&偏移&事后占用\\\midrule
234&170&申请&86&128&648&908\\
237&1716&换出&78&128&520&780\\
238&171&申请&87&128&520&908\\
240&174&申请&88&2&908&910\\
241&175&申请&89&16&910&926\\
243&1233&释放&89&16&910&910\\
244&1718&换出&88&2&908&908\\
245&1720&换出&77&128&780&780\\
246&177&申请&90&128&780&908\\\bottomrule
\end{tabular}\end{table}

这一局部包含大小为 2、16、128 的不同对象。小对象释放可以增加空闲总量，但未必立即提供适合大对象的连续段；分配过程因此在每次状态变化后重新扫描地址。在第 244、245 个位置连续换出两个对象后，大小为 128 的对象 90 被放入 $[780,908)$。这些事件显示，搬运次数与搬运量承担不同的解释功能：换出小对象增加事件数，却可能只贡献很小的数据量。

对象 78 的逻辑生命在换出后仍继续，其数据在配对换入完成后才能再次被使用。与此同时，原地址已经交给对象 87。下一问据此加入“对象 78 本段换出结束先于对象 87 申请”的边，把本节序列中的合法复用转化为时间上的等待约束。地址方案由此既决定空间可行性，也参与流水并行关系的形成。
